\section*{Hoofdstuk \ref{ch:g11:enzymes-and-phenotype} --- Enzymen en het fenotype}

\begin{solution}{exo:g11:enzymes-and-phenotype:1}
\omterm{def:g11:enzymes-and-phenotype:phenotype}{Genotype}: de verzameling gedragen \omterm{def:g10:universal-dna:gene}{allelen}. \omterm{def:g11:enzymes-and-phenotype:phenotype}{Fenotype}: de verzameling
waarneembare eigenschappen, op de moleculaire, de cellulaire en de schaal
van het organisme.
\end{solution}

\begin{solution}{exo:g11:enzymes-and-phenotype:2}
De holte van het \omterm{def:g11:enzymes-and-phenotype:enzyme}{enzym} waar het \omterm{def:g11:enzymes-and-phenotype:enzyme}{substraat} bindt en wordt omgezet.
Specificiteit: de plaats past bij één \omterm{def:g11:enzymes-and-phenotype:enzyme}{substraat} (of één familie
substraten) en katalyseert één reactie.
\end{solution}

\begin{solution}{exo:g11:enzymes-and-phenotype:3}
Rond \qty{37}{\celsius}; bij \qty{45}{\celsius} blijft ongeveer een derde over.
\end{solution}

\begin{solution}{exo:g11:enzymes-and-phenotype:4}
Moleculair: tyrosinase ontbreekt of werkt niet, er wordt geen melanine
gemaakt. Cellulair: pigmentcellen aanwezig maar leeg van melanine.
Organisme: wit haar, een heel bleke huid en bleke ogen, gevoeligheid voor
licht en zonnebrand.
\end{solution}

\begin{solution}{exo:g11:enzymes-and-phenotype:5}
De vouwing van pepsine werkt alleen in sterk zuur (pH 2); de darm is
neutraal tot licht basisch, waar de kromme geen werkzaamheid toont.
\end{solution}

\begin{solution}{exo:g11:enzymes-and-phenotype:6}
Eén werkend \omterm{def:g10:universal-dna:gene}{allel} levert de halve normale hoeveelheid \omterm{def:g11:enzymes-and-phenotype:enzyme}{enzym}, maar een
enzymmolecuul werkt duizenden keren per seconde: de halve hoeveelheid zet
nog altijd genoeg tyrosine om voor een normaal pigment. Het \omterm{def:g11:enzymes-and-phenotype:phenotype}{fenotype} is
verzadigd; alleen twee falende \omterm{def:g10:universal-dna:gene}{allelen} worden zichtbaar.
\end{solution}

\begin{solution}{exo:g11:enzymes-and-phenotype:7}
Hun albinisme komt van blokkades in twee verschillende \omterm{def:g10:universal-dna:gene}{genen} (twee
verschillende \omterm{def:g11:enzymes-and-phenotype:enzyme}{enzymen} van de route, of van het transport van melanine).
Het kind erft van elke ouder één werkend \omterm{def:g10:universal-dna:gene}{allel} van het \omterm{def:g10:universal-dna:gene}{gen} waarvan de
blokkade elders zit, dus werken beide \omterm{def:g11:enzymes-and-phenotype:enzyme}{enzymen} en wordt er pigment gemaakt.
\end{solution}

\begin{solution}{exo:g11:enzymes-and-phenotype:8}
Fenylalanine hoopt zich op; tyrosine (en wat daaruit voortkomt) ontbreekt.
Omdat melanine uit tyrosine wordt gemaakt, vermindert dat tekort het
pigment: licht haar en een lichte huid zijn bij onbehandelde
fenylketonurie gewoon.
\end{solution}

\begin{solution}{exo:g11:enzymes-and-phenotype:9}
Waar melk een basisvoedsel is, levert het verteren ervan energie en
calcium zonder ziekte, en is het persistente \omterm{def:g10:universal-dna:gene}{allel} een voordeel. Waar
volwassenen geen melk drinken, verandert het \omterm{def:g10:universal-dna:gene}{allel} niets; het is alleen
"gezonder" in de omgeving die het \omterm{def:g11:enzymes-and-phenotype:enzyme}{substraat} levert.
\end{solution}

\begin{solution}{exo:g11:enzymes-and-phenotype:10}
Voor tyrosinase: moleculair past het \omterm{def:g11:enzymes-and-phenotype:enzyme}{substraat} niet meer, geen melanine;
cellulair lege pigmentcellen; op de schaal van het organisme albinisme.
Voor lactase: lactose onverteerd; darmcellen die de suiker missen en een
dikke darm die haar vergist; intolerantie.
\end{solution}

\begin{solution}{exo:g11:enzymes-and-phenotype:11}
Haar uiteinden blijven warmer dan \qty{35}{\celsius}, het \omterm{def:g11:enzymes-and-phenotype:enzyme}{enzym} werkt
nauwelijks, en het katje is overal bleek. Buiten in de winter koelen de
oren en de voeten af, wordt het \omterm{def:g11:enzymes-and-phenotype:enzyme}{enzym} werkzaam, en komt de vacht die daar
teruggroeit donker op.
\end{solution}

\begin{solution}{exo:g11:enzymes-and-phenotype:12}
Moleculair: de helft van de hemoglobine is de sikkelvorm. Cellulair: bij
gewone zuurstofspanning blijven de \omterm{def:g10:cells-common-unit:cell}{cellen} rond; bij heel weinig zuurstof
vormen sommige een sikkel. Organisme: gezond in het gewone leven, met
risico op hoogte of bij het duiken --- en beschermd tegen malaria. De
omgeving (zuurstof, malaria) bepaalt of het \omterm{def:g11:enzymes-and-phenotype:phenotype}{fenotype} van de drager een
last of een voordeel is.
\end{solution}

\begin{solution}{exo:g11:enzymes-and-phenotype:13}
Bij \qty{41}{\celsius} houden de \omterm{def:g11:enzymes-and-phenotype:enzyme}{enzymen} nog het grootste deel van hun
werkzaamheid (de kromme zit vlak bij haar top); bij \qty{43}{\celsius}
zakt de werkzaamheid steil terwijl \omterm{def:g10:chemistry-of-life:families}{eiwitten} zich beginnen te ontvouwen,
en als dat ontvouwen doorzet, is het onomkeerbaar: \omterm{def:g11:enzymes-and-phenotype:enzyme}{enzymen} worden
vernield en niet alleen vertraagd.
\end{solution}

\begin{solution}{exo:g11:enzymes-and-phenotype:14}
Meer \omterm{def:g10:cells-common-unit:organelle}{mitochondriën} en haarvaten per \omterm{def:g10:muscles-and-joints:muscle}{spiervezel}; een groter hart met een
groter slagvolume; meer hemoglobine per liter bloed; grotere
glycogeenvoorraden. Elk daarvan is een \omterm{def:g11:enzymes-and-phenotype:phenotype}{fenotype} dat hetzelfde \omterm{def:g11:enzymes-and-phenotype:phenotype}{genotype} in
een andere omgeving voortbrengt: "erfelijke" kenmerken zijn de marge die
het \omterm{def:g11:enzymes-and-phenotype:phenotype}{genotype} toelaat, geen vaste waarde.
\end{solution}

\begin{solution}{exo:g11:enzymes-and-phenotype:15}
Het \omterm{def:g11:enzymes-and-phenotype:phenotype}{genotype} legt vast welk \omterm{def:g11:enzymes-and-phenotype:enzyme}{enzym} wordt gemaakt en hoe het zich gedraagt:
geen \omterm{def:g11:enzymes-and-phenotype:enzyme}{enzym} 1 (fenylketonurie), geen tyrosinase (albinisme), een
warmtegevoelige tyrosinase (siamees). De omgeving kan de uitkomst alleen
veranderen binnen wat het \omterm{def:g11:enzymes-and-phenotype:enzyme}{enzym} toelaat: een dieet neemt het \omterm{def:g11:enzymes-and-phenotype:enzyme}{substraat} weg
en redt de hersenen, maar kan het ontbrekende \omterm{def:g11:enzymes-and-phenotype:enzyme}{enzym} niet maken; kou laat
het siamese \omterm{def:g11:enzymes-and-phenotype:enzyme}{enzym} werken, maar geen enkele temperatuur laat dat van de
albino werken; en niets in de omgeving geeft de albino zijn pigment. Het
\omterm{def:g11:enzymes-and-phenotype:phenotype}{genotype} bepaalt de marge; de omgeving kiest daarbinnen.
\end{solution}

\begin{solution}{pb:g11:enzymes-and-phenotype:1}
\textbf{1.} Romp \qty{38}{\celsius}: ongeveer 2\%, bleek. Poten
\qty{36}{\celsius}: ongeveer 20\%, bleek. Voeten \qty{33}{\celsius}: 80\%,
donker. Oren en snuit \qty{31}{\celsius}: ruim 90\%, donker.

\textbf{2.} Bleek lijf en bleke poten, donkere voeten, oren en snuit (en
staart): het siamese patroon.

\textbf{3.} Overal 100\%: een gelijkmatig gepigmenteerde kat.

\textbf{4.} In de baarmoeder is elk deel van het katje op
\qty{38}{\celsius}, is het \omterm{def:g11:enzymes-and-phenotype:enzyme}{enzym} overal onwerkzaam, en wordt er vóór de
geboorte geen pigment afgezet; de punten worden donker naarmate de
uiteinden na de geboorte afkoelen.

\textbf{5.} Rond \qty{35.5}{\celsius}, waar de werkzaamheid door de 30\%
zakt: het steile dalende stuk van de kromme tussen 33 en
\qty{37}{\celsius} legt de grens vast.

\textbf{6.} Donker: bij \qty{30}{\celsius} is het \omterm{def:g11:enzymes-and-phenotype:enzyme}{enzym} in de groeiende
haren volop werkzaam.

\textbf{7.} Bleek: bij \qty{38}{\celsius} is het \omterm{def:g11:enzymes-and-phenotype:enzyme}{enzym} onwerkzaam.

\textbf{8.} Pigment wordt in het haar afgezet terwijl het groeit en kan er
daarna niet meer uit of bij; de kleur verandert pas wanneer de haren
uitvallen en worden vervangen door nieuwe die bij de plaatselijke
temperatuur zijn gegroeid.

\textbf{9.} Het \omterm{def:g11:enzymes-and-phenotype:phenotype}{genotype} van de \omterm{def:g10:cells-common-unit:cell}{cellen} van de plek is vóór en na gelijk,
en gelijk aan dat van de rest van de romp; alleen de temperatuur van het
\omterm{def:g11:enzymes-and-phenotype:enzyme}{enzym} veranderde, en daarmee het pigment. De ingreep werkt op de
werkzaamheid van het \omterm{def:g11:enzymes-and-phenotype:enzyme}{enzym}, een moleculaire gebeurtenis.

\textbf{10.} Alle \omterm{def:g10:cells-common-unit:cell}{cellen} van de kat stammen van één eicel af en dragen
dezelfde \omterm{def:g10:universal-dna:gene}{allelen}; en de proeven met de kaalgeschoren plek veranderen de
kleur van een streek zonder haar \omterm{def:g10:cells-common-unit:cell}{cellen} te veranderen --- een \omterm{def:g10:universal-dna:gene}{allel} valt
niet met een koelelement om te zetten.

\textbf{11.} Moleculair: een tyrosinase die alleen beneden ongeveer
\qty{35}{\celsius} werkt. Cellulair: pigmentcellen vol melanine in de
uiteinden, leeg in de romp. Organisme: een bleke kat met donkere punten.

\textbf{12.} Het \omterm{def:g10:chemistry-of-life:families}{aminozuur} zit op de plaats waar het de vouwing
stabiliseert; de vervanger houdt de vouwing bij lage temperatuur bijeen
maar laat haar enkele graden hoger loszitten, zodat de \omterm{def:g11:enzymes-and-phenotype:enzyme}{actieve plaats} haar
vorm eerder verliest dan bij het gewone \omterm{def:g11:enzymes-and-phenotype:enzyme}{enzym}.

\textbf{13.} Het \omterm{def:g11:enzymes-and-phenotype:enzyme}{enzym} van het gewone \omterm{def:g10:universal-dna:gene}{allel} is bij lichaamstemperatuur
volop werkzaam; de halve normale hoeveelheid is nog altijd ruim genoeg om
overal pigment te maken. Het siamese \omterm{def:g10:universal-dna:gene}{allel} is er recessief tegenover.

\textbf{14.} Moleculair: bij geen enkele temperatuur een werkend \omterm{def:g11:enzymes-and-phenotype:enzyme}{enzym},
tegenover een \omterm{def:g11:enzymes-and-phenotype:enzyme}{enzym} dat onder voorwaarden werkt. Cellulair: nergens
melanine, tegenover melanine in de koele streken. Organisme: een geheel
witte kat met roze ogen, tegenover het puntenpatroon. Voor de albino
verandert kou niets.

\textbf{15.} Dat zij een temperatuurgevoelig \omterm{def:g10:universal-dna:gene}{allel} van hetzelfde \omterm{def:g10:universal-dna:gene}{gen}
dragen; waarschijnlijk een vergelijkbare substitutie die het \omterm{def:g11:enzymes-and-phenotype:enzyme}{enzym}
minder stabiel maakt, in elke \omterm{def:g10:biodiversity-scales:species}{soort} afzonderlijk ontstaan.

\textbf{16.} Twintigvoudig; \omterm{def:g11:enzymes-and-phenotype:enzyme}{enzym} 1, dat fenylalanine in tyrosine omzet.

\textbf{17.} Overtollige fenylalanine hindert juist de \omterm{def:g10:cells-common-unit:cell}{cellen} van de zich
ontwikkelende hersenen (hun aanvoer van andere \omterm{def:g10:chemistry-of-life:families}{aminozuren} en hun
rijping); huidcellen verdragen haar. Het moleculaire gebrek is overal
hetzelfde; het cellulaire gevolg niet.

\textbf{18.} Een vijfde van \num{1200} is \qty{240}{\micro mol/L}, onder
\num{360}: veilig. Zonder het \omterm{def:g11:enzymes-and-phenotype:enzyme}{enzym} volgt het bloedgehalte eenvoudigweg
de inname, zodat het \omterm{def:g11:enzymes-and-phenotype:phenotype}{fenotype} door de aanvoer van \omterm{def:g11:enzymes-and-phenotype:enzyme}{substraat} --- de
omgeving --- wordt bepaald en niet door het ontbrekende \omterm{def:g11:enzymes-and-phenotype:enzyme}{enzym} alleen.

\textbf{19.} Fenylalanine is nodig om de \omterm{def:g10:chemistry-of-life:families}{eiwitten} van het lichaam te
bouwen en kan niet worden gemaakt; zonder enige fenylalanine stopt de
groei. De inname moet genoeg zijn voor de eiwitsynthese en niet meer, en
dat vergt meten en bijstellen.

\textbf{20.} Het \omterm{def:g11:enzymes-and-phenotype:phenotype}{genotype} legde het \omterm{def:g11:enzymes-and-phenotype:enzyme}{enzym} vast (warmtegevoelig bij de kat,
afwezig bij het kind); de omgeving bepaalde de voorwaarden ervoor
(huidtemperatuur; fenylalanine in de voeding); het \omterm{def:g11:enzymes-and-phenotype:phenotype}{fenotype} is het
resultaat van het \omterm{def:g11:enzymes-and-phenotype:phenotype}{genotype} zoals het in een gegeven omgeving tot uiting
komt.
\end{solution}
